\begin{theorem}[Every Element Is Either $1$ or a Successor]
\label{thm:every-element-is-one-or-a-successor}
Let $(P,S,1)$ be a Peano system. For every $x \in P$, either $x=1$ or there
exists $u \in P$ such that $S(u)=x$.
\hyperref[prf:every-element-is-one-or-a-successor]{Go to proof.}
\end{theorem}
\end{theorembox}
\LeanFormalizes{thm:every-element-is-one-or-a-successor}{lra-lean}{LRA.VolumeII.PeanoSystems.BasicTheorems}{EveryElementIsOneOrASuccessor}{pending}

\begin{formalizationrecord}{Lean 4 Verification Record}
\FormalizationSource{Landau, \emph{Foundations of Analysis}, Section 1; Feferman, \emph{The Number Systems}, Theorem 3.3}
\Formalizes{thm:every-element-is-one-or-a-successor}
\textbf{Lean module:} \texttt{LRA.VolumeII.PeanoSystems.BasicTheorems}.\par
\LeanDeclaration{every\_element\_is\_one\_or\_successor}
\textbf{Status:} checked against \texttt{lra-lean}.\par

\begin{leancode}
theorem every_element_is_one_or_successor
    (ps : PeanoSystem) :
    forall element : ps.carrier,
      Or (element = ps.one)
        (Exists (fun predecessor : ps.carrier =>
          ps.successor predecessor = element)) := by
  let D : ps.carrier -> Prop :=
    fun candidate_element =>
      Or (candidate_element = ps.one)
        (Exists (fun predecessor : ps.carrier =>
          ps.successor predecessor = candidate_element))

  apply induction_principle ps D

  case base_case =>
    exact Or.inl rfl

  case successor_step =>
    intro element induction_hypothesis
    exact Or.inr (Exists.intro element rfl)
\end{leancode}
\end{formalizationrecord}

\begin{remark*}[Standard quantified statement]
For a Peano system $(P,S,1)$,
\[
\forall x \in P\;\bigl(x=1 \lor \exists u \in P\;(S(u)=x)\bigr).
\]
\end{remark*}

\begin{remark*}[Predicate reading]
\[
\operatorname{SuccessorCoversNonbase}(S,1,P).
\]
\end{remark*}

\begin{remark*}[Interpretation]
Every element of a Peano system is reached by starting at the distinguished
first element and moving forward by successor. The only possible exception to
having a predecessor is the element $1$ itself.
\end{remark*}

\begin{remark*}[Historical note]
This theorem is one of the basic successor consequences of Peano's
axiomatization \citep{PeanoArithmeticesPrincipia1889}. In the present
organization it corresponds directly to Theorem 3.3 in
\citet{FefermanNumberSystems1964}. Feferman uses this result to show that
the distinguished initial element is the only possible non-successor in a
Peano system.
\end{remark*}

\begin{dependencies}
\begin{itemize}
  \item \hyperref[def:peano-system]{Peano System}
  \item \hyperref[ax:peano-base-in-set]{Distinguished Element Lies in the Underlying Set}
  \item \hyperref[ax:peano-successor-closure]{Closure Under Successor}
  \item \hyperref[thm:induction-principle-for-peano-system]{Induction Principle for a Peano System}
\end{itemize}
\end{dependencies}

